\begin{figure*}[t]
\centering
\scriptsize
\setlength{\fboxsep}{1em}
\setlength{\fboxrule}{1pt}
\fcolorbox{black}{white}{
\parbox{0.955\textwidth}{
\begin{alltt}


You will be provided instructions for completing a task followed by a task to complete.

Instructions:

1. Understand the structure of the dataset: The dataset consists of news articles and their corresponding one-sentence summaries. The 'document' feature represents the input news article, and the 'summary' feature represents the one-sentence summary of the article.

2. Familiarize yourself with the purpose of the dataset: The goal is to create a short, one-sentence summary that answers the question “What is the article about?”.

3. Read the entire article carefully: Make sure to understand the main points of the article. The articles cover a wide variety of domains such as News, Politics, Sports, Weather, Business, Technology, Science, Health, Family, Education, Entertainment, and Arts.

4. Generate a one-sentence summary: After understanding the article, generate a one-sentence summary that captures the main point of the article. The summary should be concise and clear, and it should accurately represent the content of the article.

5. Check your summary: Make sure your summary accurately reflects the main point of the article. It should not include any personal opinions or interpretations, and it should not introduce any new information that was not present in the article.

6. Practice with different articles: The more you practice, the better you will become at generating accurate and concise summaries. Try to practice with articles from different domains to improve your summarization skills.

Remember, the goal is not to generate the longest or most detailed summary, but to capture the main point of the article in a single sentence.

\#\#\#

Article: Head coach Stuart Lancaster's World Cup preparations suffered a blow as for the first 70 minutes a largely first-choice XV struggled to deal with French power. Two late tries flattered the visitors, who have one game left before launching their World Cup campaign against Fiji on 18 September. "We gave away penalties and our discipline was shocking," said Robshaw. "Whether it was rust, or nerves, it wasn't good enough. Credit to France, they put us under pressure and made us make mistakes. "We gave too many penalties away, but in the second half we came out and played well but couldn't quite get over the line in the end," he told Sky Sports. "We can't give teams like France and other quality sides head starts like we did. "We'll look long and hard at ourselves, because we let ourselves down in the first half. We played well in phases but you can't wait for 40 minutes to do that." Late tries from Danny Cipriani and Jonathan Joseph made it close on the scoreboard but Lancaster was left with much to ponder by a disappointing team display in the Stade de France. Media playback is not supported on this device The head coach, who must announce his final squad by 31 August, added: "We've got to get our discipline at the breakdown - we can't give France easy position like we did. We need to improve on that, because all the little mistakes add up. "The bench made a difference. It upped the energy levels and we scored some good tries. I won't gloss over what went on before that, because it was too little too late. "There are a few players who have given me food for thought, those guys who came on and gave us the energy we needed and made a difference. "I need to have a sleep on this game and think about my final squad. We're two weeks away from playing Ireland and four weeks away from Fiji in the World Cup and we'll expect a reaction next time." England host Ireland in their final World Cup warm-up game at Twickenham on 5 September. Fly-half Ford told Sky Sports: "I thought we might snatch that at the end but we had hardly any ball in the first half and gave away too many penalties. We played some good stuff in the second-half. "In the first half a lot of our undoing was down to ourselves. We just weren't good enough in that first half and there's no excuse for that. We let them build up the score and that made it hard for us. "It was frustrating and we had to think of ways to adapt and that was constantly going through our minds. We tried to get ourselves out of a hole. "We've got to turn up against Ireland now and make sure that we win. Our basics have got to be world class."

Follow the instructions to generate an explanation that reasons towards the correct answer to the task above. End the explanation with the correct answer. [/INST] Explanation: 

The article discusses England's rugby team's preparation for the World Cup, specifically their match against France. The team's performance was disappointing, with poor discipline and mistakes leading to a loss. The article highlights the struggles of the team's first-choice XV and the impact of late tries from Danny Cipriani and Jonathan Joseph. The head coach, Stuart Lancaster, expresses his disappointment and emphasizes the need for improvement in discipline and energy levels. The article also mentions the team's upcoming match against Ireland and their final World Cup warm-up game.

Correct answer: The article is about England's rugby team's preparation for the World Cup and their disappointing performance against France.

Therefore, the answer to the task is below. Give the answer in the shortest form possible that will still be correct.

Summarize the above article in 1 sentence.

England's rugby team struggled in their World Cup preparation match against France, with poor discipline and mistakes leading to a loss, and head coach Stuart Lancaster emphasized the need for improvement.
\end{alltt}
}}
\caption{A complete running example of XSUM.}
\label{fig:summarization_xsum_example}
\end{figure*}
